% Root and B full literal audits passed; no numerical premise.
\subsection{Four cross states opposite the small nonprojective corners}

\begin{lemma}[The minimal right-zero graph surjects on fundamental groupoids]
\label{C:right-zero-graph-fundamental-quotient}
Let a finite monoid $N$ have minimal ideal $K=\{p,q\}$, a
right-zero band, and let $S$ be an $(N,M)$-biset.
The actual inclusion
\[
 (K^1\ltimes S)\longrightarrow(N\ltimes S)
\]
induces a bijection on components and a surjection on every
fundamental-groupoid hom set. These assertions and the canonical
maps from action categories to their fundamental groupoids are
natural for the commuting right $M$-action.
\end{lemma}
\begin{proof}
For any arrow $n_s:s\to ns$ of the larger action category,
composition with $p_{ns}:ns\to pns$ has label $pn\in K$.
After inverting arrows,
\[
 n_s=p_{ns}^{-1}(pn)_s.
\]
Both arrows on the right belong to the smaller action category.
Thus each new arrow joins objects already joined there, and
every path in the larger fundamental groupoid is the image
of a path in the smaller one. The identity-on-objects inclusion
and the defining equations commute with every right actor.
Groupoid completion itself is functorial, giving all the stated
naturality.
\end{proof}

\begin{lemma}[The low-degree component and groupoid comparisons]
\label{C:action-space-low-degree-comparison}
Let a finite monoid $M$ act on a CW space $Z$, with the usual
natural action-bar projection to $BM$.
If every component of $Z$ has zero first rational homology,
projection to $T=\pi_0Z$ gives a surjection
\[
 H^2\bigl(\operatorname{hocofib}(T_{hM}\to BM);\mathbb Q\bigr)
 \longrightarrow
 H^2\bigl(\operatorname{hocofib}(Z_{hM}\to BM);\mathbb Q\bigr).
\]
If $Z$ is connected, then
\[
 H_2\bigl(\operatorname{hocofib}(Z_{hM}\to BM);\mathbb Q\bigr)
   =H_1(Z;\mathbb Q)\otimes_{\mathbb QM}\mathbb Q.
\]
Consequently a natural map of connected coefficient spaces
inducing an isomorphism on $H_1$ induces an isomorphism on the
second homology of these actual cofibers. In particular this
applies to the classifying map from an action category to its
fundamental groupoid.
\end{lemma}
\begin{proof}
Use the augmented singular or simplicial chain complex and
tensor it with the free left monoid bar resolution. Its
suspension is the reduced cofiber chain model. The kernel of
the projection onto $\widetilde{\mathbb QT}$ has zero
degree-zero homology, with first homology $H_1(Z;\mathbb Q)$.
The first-quadrant total complex therefore gives the displayed
surjection when that first homology vanishes.
For connected $Z$, the only term contributing to total degree
two after suspension is its first homology in bar degree zero;
there is no lower nonzero coefficient row from which a
differential could enter. This gives the tensor formula and
its naturality. These statements require neither bounded
homology nor finite CW type for $Z$.
\end{proof}

\begin{theorem}[Four-state reduction for an arbitrary two-element right-zero ideal]
\label{C:right-zero-minimal-two-four-states}
Let a finite monoid $N$ have minimal ideal a two-element
right-zero band. For every $(N,M)$-biset $S$ with $|S|\leq4$,
a collage of type $\mathbb CP^2$ yields an ordinary right
$M$-act cone on at most $|S|$ states of that actual type.
For each such fixed $N$, the least arrow count among these
collages realizing $\mathbb CP^2$ is exactly $21+|N|$.
\end{theorem}
\begin{proof}
Write $P,Q$ for the two minimal retractions, with common
image $U$. All $N$-actors fix $U$ pointwise, and $BN$ is
contractible. The cases with at most three states were proved
in Lemma~\ref{C:right-zero-minimal-two-three-states}.
A singleton $U$ gives a common fixed state; $U=S$ gives an
actual ordinary act cone.

If $P=Q$, this common map is a biset retraction onto $U$:
$nP=P$ and $Pn=P$ follow respectively from $nK=K$
pointwise and $Kn\subseteq K$. The core collage $Y$ is an
ordinary $U$-act cone and a homotopy retract of the original
collage $X$. The relative quotient is a suspension.
If $H^2(Y;\mathbb Q)=0$, the projective-plane square is
therefore impossible. Otherwise the split integral homology
argument used in Theorem~\ref{C:right-zero-five-state-reduction}
makes $Y\to X$ an integral homology equivalence: degree two
is an isomorphism by splitting, degree four is nonzero by the
square and hence an isomorphism by splitting, and all other
reduced homology vanishes. Both spaces are simply connected
because $Y$ is a retract of $X$. Homological Whitehead gives
the required actual equivalence.

Suppose now that $P\ne Q$, and use the natural graph $\Gamma$
of the $K^1$ action. If $|U|=3$, its single edge is nonloop,
so its two components are trees. By
Lemma~\ref{C:right-zero-graph-fundamental-quotient}, the space
$Z=B(N\ltimes S)$ has the same two components and each is
simply connected. Its component set $T$ is a two-state right
act. Lemma~\ref{C:action-space-low-degree-comparison} makes
the second cohomology of its ordinary act cone surject onto
that of the collage. The finite-line lemma kills every square.

The remaining graph has two vertices and two edges and is
connected. Its fundamental group is infinite cyclic. The
same groupoid-surjection lemma says that $Z$ is connected
and its fundamental group is a quotient of $\mathbb Z$.
If this quotient is finite, $H_1(Z;\mathbb Q)=0$ and the
collage has zero second rational homology. If it is infinite,
the actual composite
\[
 B(K^1\ltimes S)\longrightarrow B(N\ltimes S)
       \longrightarrow B\Pi(N\ltimes S)
\]
is a weak equivalence: its domain is the unicyclic graph and
its target is a connected $K(\mathbb Z,1)$, and the induced
surjection $\mathbb Z\to\mathbb Z$ is an isomorphism.
This entire comparison is $M$-equivariant. Taking bars gives
an equivalence on cofibers between the graph cofiber and the
fundamental-groupoid cofiber. Meanwhile the actual map from
the original cofiber to the latter induces an isomorphism on
$H_2(-;\mathbb Q)$ by
Lemma~\ref{C:action-space-low-degree-comparison}.

A nonzero second homology forces every right actor to have
degree $+1$ on the graph cycle. The proof of
Lemma~\ref{C:unicyclic-graph-correspondence} makes every
positive rational product on its cofiber zero. Pullback
along the actual second-homology isomorphism therefore kills
the projective-plane square on the original collage. This
excludes the last case.

The ordinary-act reduction and its exact minimum twenty-two
give the lower bound $21+|N|$. For equality let $N$ act
trivially on the three states of the accepted eighteen-element
source model. Since $BN$ is contractible, its collage is
actually equivalent to the twenty-two-arrow act cone and
has $18+3+|N|$ arrows.
\end{proof}

\begin{theorem}[Four states opposite the three-element right-zero band]
\label{C:right-zero-three-four-states}
Let $N$ be the identity-adjoined three-element right-zero band,
and let $|S|\leq4$. A collage of type $\mathbb CP^2$ yields
an ordinary right $M$-act cone on at most $|S|$ states of that
actual type. In particular this fixed opposite corner has
exact least arrow count twenty-five in the stated state range.
\end{theorem}
\begin{proof}
Use the natural graph with core vertices $U$, the outside
states as additional vertices, and one edge from each outside
state for each of the three right-zero labels. Its core is
fixed pointwise by $N$. A singleton core gives a common fixed
state; a trivial action gives an ordinary act. If the three
retractions coincide, their common biset retraction reduces
a projective plane to an actual ordinary act by the same
split-homology argument as above.

With three core states and one outside state $z$, three
distinct endpoints give a tree, hence a contractible collage.
Two distinct endpoints give a unicyclic component and one
isolated core vertex. If any right actor sends $z$ into the
core, it kills the cycle class, so its coinvariants vanish.
The component comparison of
Lemma~\ref{C:graph-component-eigenvalues} makes every degree-two
class pull back from a two-state act cone, and its square is
zero. Otherwise every right actor fixes $z$ and both of its
core endpoints; either endpoint is a common fixed state.

It remains that $U=\{u_0,u_1\}$ and there are two outside
states $z,w$. Record the three retraction values at an outside
state as a word $v_s\in\{0,1\}^3$. At least one of the two
words is nonconstant, or all minimal retractions coincide.
The graph is therefore connected and has first Betti number
three. If a right actor induces the vertex map
$\sigma:\{0,1\}\to\{0,1\}$, commutation says that the word
of the image of $s$ is $\sigma\circ v_s$; core states have
words $000$ and $111$.

Suppose first that both outside words are nonconstant.
A constant vertex map then sends every state to that constant
core vertex, and kills the entire first homology. It would
make the collage's second homology zero, so such actors may
be excluded in the putative projective-plane case.
If every vertex map is the identity, there is a common fixed
core state. Otherwise a vertex swap occurs, which forces
$v_w$ to be the complement of $v_z$. The words are distinct,
so an actor fixing the core fixes all states, while an actor
swapping it is the unique simultaneous interchange of
$u_0,u_1$ and $z,w$. The graph action image is exactly $C_2$.
This involution acts freely on its four vertices and six
edges. Its quotient graph has two vertices and three edges,
so the rational first-homology coinvariants have dimension
two. The connected graph-Tor comparison gives
$\dim H_2(X;\mathbb Q)=2$, again excluding a projective plane.

Finally suppose, after relabeling, that $v_z$ is nonconstant
and $v_w=000$. A vertex swap is impossible, since the
complement of $v_z$ is not the word of either outside state.
A constant-one vertex map collapses all states to $u_1$ and
kills first homology. Unless the second homology is zero,
all vertex maps are therefore the identity or constant zero;
they all fix $u_0$. The common-fixed-state argument kills
the degree-two products. These cases are exhaustive.
The lower and upper twenty-five bounds follow from the
ordinary-act reduction and the trivial-action construction.
\end{proof}

\begin{corollary}[The nonprojective order-four opposite corners]
\label{C:nonprojective-four-corners-four-states}
For every four-element monoid $N$ whose trivial right rational
module is nonprojective, a collage with at most four cross
states and type $\mathbb CP^2$ has at least twenty-five arrows.
The bound is attained for each of the six such monoids.
In particular the seventeen-arrow profile $(9,4,4)$ is
excluded, including all nonprojective opposite corners.
\end{corollary}
\begin{proof}
The elementary inventory in
Theorem~\ref{C:small-opposite-corner-three-states} leaves
precisely a two-element or a three-element right-zero minimal
ideal. Apply the two preceding theorems. The projective
opposite corners in the stated seventeen-arrow profile were
already excluded by
Corollary~\ref{C:projective-opposite-corner-twenty-one}.
\end{proof}
